\begin{description}
\item[(a)] Energy of a photon:
  \[ E = hf \]
  while the momentum:
  \[ \vec{p} = \frac{hf}{c}\vec{n} \]
  \hfill(2.0 pt)

  where $\vec{n}$ is an unit vector in the direction of the motion of the
  photon. Substituting this into energy-momentum transformation law we have:
  \begin{align*}
    \frac{hf_L}{c} &= \gamma\left(\frac{hf_R}{c} + p_{x_R}\frac{v}{c}\right)\\
    p_{x_L} &= \gamma\left(p_{x_R} + \frac{hf_R v}{c^2}\right)\\
    p_{y_L} &= p_{y_R}\\
    p_{z_L} &= p_{z_R}
  \end{align*}
  \hfill(2.0 pt)

  in our case:
  \begin{align*}
    p_{x_L} &= |p_L|\cos\theta_L, & p_{y_L} &= |p_L|\sin\theta_L,
      & p_{z_L} &= 0\\
    p_{x_R} &= |p_R|\cos\theta_R, & p_{y_R} &= |p_R|\sin\theta_R,
      & p_{z_R} &= 0
  \end{align*}
  \hfill(3.0 pt)

  where
  \[ |p| = \frac{hf}{c} \]
  Thus,
  \begin{align*}
    f_L &= \gamma\left(f_R + \frac{v}{h}|p_R|\cos\theta_R\right)\\
        &= \gamma\left(f_R + \frac{v}{h}\,\frac{hf_R}{c}\cos\theta_R\right)\\
    \therefore\ f_L &= \gamma f_R\left(1 + \frac{v}{c}\cos\theta_R\right)
  \end{align*}
  \hfill(2.0 pt)
  \begin{align*}
    \cos(\theta_L) &= \frac{p_{x_L}}{|p_L|} = \frac{c\,p_{x_L}}{hf_L}
      = \frac{c\gamma\left(p_{x_R} + \dfrac{hf_R v}{c^2}\right)}
             {h\gamma f_R\left(1 + \dfrac{v}{c}\cos\theta_R\right)}\\
    &= \frac{c\left(\dfrac{hf_R}{c}\cos\theta_R + \dfrac{hf_R v}{c^2}\right)}
            {hf_R\left(1 + \dfrac{v}{c}\cos\theta_R\right)}\\
    \therefore\ \cos\theta_L
      &= \frac{\cos\theta_R + \dfrac{v}{c}}
              {1 + \dfrac{v}{c}\cos\theta_R}
  \end{align*}
  \hfill(2.0 pt)
\item[(b)] Now just by plugging in the values of cosine we see that in case
  (i) (orange arrow) $\cos\theta_L = 1$, So the photon keeps moving in the
  same direction, in case (ii) (Green arrow) $\cos\theta_L = 0$, in case
  (iii) (Yellow arrow) $\cos\theta_L = v/c$, (iv) (Grey arrow)
  $\cos\theta_L = -1$.
  % FIGURE NEEDED: two XY-axis diagrams side by side ("Lab Frame" and "Rest
  % Frame") showing the four colour-coded direction arrows (orange, green,
  % yellow, grey) of the beam in each frame; 2022 theory solutions, page 19.
\item[(c)] From the figure we see that:
  \[ KP = OP - OK\sin\theta = R - r\sin\theta \]
  \hfill(1.0 pt)

  where $\theta = \omega t$.
  % FIGURE NEEDED: geometry sketch of the circular orbit: centre O, particle
  % K on the circle, angle theta between OK and the OX axis marked at O and
  % at K, and two parallel rays from O and K towards the distant observer P;
  % 2022 theory solutions, page 19.

  Let $t_L = T$ be the time at which photon reaches $P$ after leaving $K$ at
  a time $t$ (in lab frame) then:
  \begin{align*}
    c(T - t) &= KP = R - r\sin\theta \approx R\\
    \therefore\ t &= t_L - \frac{R}{c}
  \end{align*}
  \hfill(1.0 pt)

  This photon makes angle $\theta = \omega t$ to the direction of its motion
  in lab frame. Now,
  \begin{align*}
    \cos\theta_L &= \frac{\cos\theta_R + \dfrac{v}{c}}
                        {1 + \dfrac{v}{c}\cos\theta_R}\\
    \cos\theta_L + \frac{v}{c}\cos\theta_R\cos\theta_L
      &= \cos\theta_R + \frac{v}{c}\\
    \cos\theta_R\left(1 - \frac{v}{c}\cos\theta_L\right)
      &= \cos\theta_L - \frac{v}{c}\\
    \therefore\ \cos\theta_R
      &= \frac{\cos\theta_L - \dfrac{v}{c}}
              {1 - \dfrac{v}{c}\cos\theta_L}
  \end{align*}
  \hfill(2.0 pt)
  \begin{align*}
    f_L &= \gamma f_R\left(1 + \frac{v}{c}\cos\theta_R\right)\\
    &= \gamma f_R\left(1 + \frac{\dfrac{v}{c}
        \left(\cos\theta_L - \dfrac{v}{c}\right)}
        {1 - \dfrac{v}{c}\cos\theta_L}\right)\\
    &= \gamma f_R\left(\frac{1 - \dfrac{v}{c}\cos\theta_L
        + \dfrac{v}{c}\cos\theta_L - \left(\dfrac{v}{c}\right)^2}
        {1 - \dfrac{v}{c}\cos\theta_L}\right)\\
    &= \frac{f_R}{\gamma\left(1 - \dfrac{v}{c}\cos\theta_L\right)}
     = \frac{f_R}{\gamma\left(1 - \dfrac{v}{c}\cos(\omega t)\right)}\\
    \therefore\ f_L
      &= \frac{f_R}{\gamma\left(1 - \dfrac{v}{c}
         \cos\left[\omega\left(t_L - \dfrac{R}{c}\right)\right]\right)}
  \end{align*}
  \hfill(3.0 pt)
\item[(d)] \hspace{0pt}
  \begin{align*}
    \Delta\Omega_R &= -\Delta(\cos\theta_R)\cdot\Delta\phi\\
    &= \left[\cos\theta_R - \cos(\theta_R + \Delta\theta_R)\right]
       \cdot\Delta\phi\\
    &= \left[\cos\theta_R - \cos\theta_R\cos(\Delta\theta_R)
       + \sin\theta_R\sin(\Delta\theta_R)\right]\cdot\Delta\phi\\
    &= \left[\cos\theta_R - \cos\theta_R\cdot(1)
       + \sin\theta_R(\Delta\theta_R)\right]\cdot\Delta\phi\\
    \Delta\Omega_R &= \sin\theta_R(\Delta\theta_R)(\Delta\phi)\\
    \therefore\ \Delta\Omega_L &= \sin\theta_L(\Delta\theta_L)(\Delta\phi)
  \end{align*}
  \hfill(3.0 pt)
  \begin{align*}
    \sin^2\theta_L &= 1 - \cos^2\theta_L\\
    &= 1 - \left(\frac{\cos\theta_R + \dfrac{v}{c}}
        {1 + \dfrac{v}{c}\cos\theta_R}\right)^2\\
    &= \frac{1 + 2\dfrac{v}{c}\cos\theta_R
        + \left(\dfrac{v}{c}\right)^2\cos^2\theta_R - \cos^2\theta_R
        - 2\dfrac{v}{c}\cos\theta_R - \left(\dfrac{v}{c}\right)^2}
        {\left(1 + \dfrac{v}{c}\cos\theta_R\right)^2}\\
    &= \frac{\left(1 - \left(\dfrac{v}{c}\right)^2\right)
        \left(1 - \cos^2\theta_R\right)}
        {\left(1 + \dfrac{v}{c}\cos\theta_R\right)^2}
     = \frac{\sin^2\theta_R}
        {\gamma^2\left(1 + \dfrac{v}{c}\cos\theta_R\right)^2}\\
    \sin\theta_L &= \frac{\sin\theta_R}
        {\gamma\left(1 + \dfrac{v}{c}\cos\theta_R\right)}
  \end{align*}
  \hfill(3.0 pt)
  \begin{align*}
    \text{Also, } \Delta\theta_L \approx \sin(\Delta\theta_L)
      &= \frac{\sin(\Delta\theta_R)}
         {\gamma\left(1 + \dfrac{v}{c}\cos\theta_R\right)}\\
    \Delta\theta_L &\approx \frac{\Delta\theta_R}
      {\gamma\left(1 + \dfrac{v}{c}\cos\theta_R\right)}
  \end{align*}
  \hfill(2.0 pt)
  \begin{align*}
    \frac{\Delta\Omega_L}{\Delta\Omega_R}
      &= \frac{\sin\theta_L(\Delta\theta_L)(\Delta\phi)}
              {\sin\theta_R(\Delta\theta_R)(\Delta\phi)}\\
      &= \frac{\sin\theta_R(\Delta\theta_R)}
              {\gamma^2\left(1 + \dfrac{v}{c}\cos\theta_R\right)^2}
         \cdot\frac{1}{\sin\theta_R(\Delta\theta_R)}\\
    \therefore\ \Delta\Omega_L
      &= \frac{\Delta\Omega_R}
              {\gamma^2\left(1 + \dfrac{v}{c}\cos\theta_R\right)^2}
  \end{align*}
  \hfill(2.0 pt)
\item[(e)] Let $N(\theta_L)\Delta\Omega_L\Delta T$ be the number of photons
  arriving in the vicinity of $P$ within the element of solid angle
  $\Delta\Omega_L$ and time interval $T, T + \Delta T$ (it is important that
  $\Delta T \neq \Delta t$ where $\Delta t$ is the emission time of the same
  photons (in lab frame)). These being photons emitted by K into the solid
  angle $\Delta\Omega$ in the time interval $t_{0R}, t_{0R} + \Delta t_{0R}$,
  so that:
  \[ N(\theta_L)\Delta\Omega_L\Delta T = N\Delta\Omega_R\Delta t_{0R} \]
  \hfill(2.0 pt)

  From an earlier part:
  \begin{align*}
    cT &= ct + R - r\sin\theta_L = ct + R - r\sin(\omega t)\\
    % sic: the official solution writes sin(theta_L) here; from part (c) the
    % angle is theta = omega*t, not theta_L.
    \therefore\ c\Delta T &= c\Delta t
      - r\left[\sin(\omega t)\cos(\omega\Delta t)
      + \sin(\omega\Delta t)\cos(\omega t) - \sin(\omega t)\right]\\
    c\Delta T &= c\Delta t - r\omega\cos(\omega t)\Delta t
  \end{align*}
  \hfill(2.0 pt)
  \begin{align*}
    \frac{\Delta T}{\Delta t_{0R}}
      &= \frac{\Delta t}{\Delta t_{0R}}\frac{\Delta T}{\Delta t}
       = \gamma\left(1 - \frac{v}{c}\cos(\omega t)\right)\\
    \frac{\Delta t_{0R}}{\Delta T}
      &= \frac{1}{\gamma\left(1 - \dfrac{v}{c}\cos(\omega t)\right)}
  \end{align*}
  \hfill(3.0 pt)
  \begin{align*}
    \frac{N(t_L)}{N_R}
      &= \frac{\Delta\Omega_R\Delta t_{0R}}{\Delta\Omega_L\Delta T}\\
      &= \frac{\gamma^2\left(1 + \dfrac{v}{c}\cos\theta_R\right)^2}
              {\gamma\left(1 - \dfrac{v}{c}\cos(\omega t)\right)}\\
      &= \frac{\gamma\left(1 + \left(\dfrac{v}{c}\right)
          \dfrac{\cos\theta_L - \dfrac{v}{c}}
                {1 - \dfrac{v}{c}\cos\theta_L}\right)^2}
              {\left(1 - \dfrac{v}{c}\cos\theta_L\right)}\\
      &= \frac{\gamma}
              {\gamma^4\left(1 - \dfrac{v}{c}\cos\theta_L\right)^3}\\
    N(t_L) &= \frac{N_R}
      {\gamma^3\left(1 - \dfrac{v}{c}\cos\left(t_L - R/c\right)\right)^3}
    % sic: the official solution omits the factor omega inside the cosine;
    % it should read cos[omega(t_L - R/c)].
  \end{align*}
  \hfill(3.0 pt)

  In terms of energy fluxes
  \begin{align*}
    F_0 &= \frac{hf_R N_R}{R^2} = \frac{L}{4\pi R^2}\\
    F(t_L) &= \frac{hf_L N(t_L)}{R^2}
  \end{align*}
  \hfill(2.0 pt)
  \begin{align*}
    \frac{F(t_L)}{F_0} &= \frac{hf_L N(t_L)}{hf_R N_R}\\
      &= \frac{1}{\gamma^4\left(1 - \dfrac{v}{c}
         \cos\left[\omega\left(t - \dfrac{R}{c}\right)\right]\right)^4}\\
    F(t_L) &= \frac{L}{4\pi R^2\gamma^4\left(1 - \dfrac{v}{c}
      \cos\left[\omega\left(t - \dfrac{R}{c}\right)\right]\right)^4}
  \end{align*}
  \hfill(3.0 pt)

  The radiation is strongly beamed in the direction of motion of the source
  so that a remote observer in or near the orbital plane of the source sees
  strongly pulsed radiation.
\item[(f)] The amplification (for a given $v$) is highest when
  $\cos\theta_L = 1$. \hfill(1.0 pt)
  \begin{align*}
    F(t_L) &= \frac{F_0}{\left[\gamma\left(1 - \dfrac{v}{c}\right)\right]^4}\\
    A_{\max} &= \frac{1}{\left[0.0975 \times 0.05\right]^4}\\
    A_{\max} &\approx 1.8 \times 10^9
    % sic: the official solution substitutes 0.0975 = 1/gamma^2 where the
    % formula above calls for gamma = 3.2; evaluating [gamma(1-v/c)]^{-4}
    % directly gives ~1.5e3, not 1.8e9 (and similarly for A_min below).
  \end{align*}
  \hfill(1.0 pt)
  \begin{align*}
    \text{Similarly, } A_{\min}
      &= \frac{1}{\left[0.0975 \times 1.95\right]^4}\\
    A_{\min} &\approx 770
  \end{align*}
  \hfill(1.0 pt)
\end{description}
